% !TEX root = ../main.tex
\section{Methods}

\hypertarget{case-study-the-umpqua-national-forest-forplan-model}{%
\subsection{Case study: the Umpqua National Forest FORPLAN model}\label{case-study-the-umpqua-national-forest-forplan-model}}

The case study reproduces the Umpqua National Forest FORPLAN planning model \citep{usda1987umpqua}, reconstructed in the open-source \texttt{ws3} wood-supply model \citep{paradis2026ws3}. The forest is stratified into four ecoclasses---CH-CW, CD-CP, CR-CF, and CM-CE, dominated respectively by Douglas-fir, hemlock, fir, and hemlock types---arrayed in decreasing growth potential. Seven silvicultural prescriptions are available, each with ecoclass-specific rotation-age ranges. The planning horizon is 150 years, divided into 15 periods of 10 years. The objective is maximization of net present value at a 4\% discount rate (varied in the experiments), with delivered-log prices escalating at 1\% per year for the first 50 years and constant thereafter. A harvest-flow (even-flow) constraint links total harvest volume between consecutive periods, and volumes are expressed in thousands of cubic feet (MCF). The initial forest is specified as one of a set of landbase conditions; the landbase of central interest (landbase 1) is an all-mature forest, the extreme of the disequilibrium structures the thesis studies.

\hypertarget{case-study-data}{%
\subsection{Case-study data}\label{case-study-data}}

The case study is built from the real Umpqua FORPLAN data, obtained from the public-domain Umpqua National Forest planning documents on HathiTrust. The per-age yield curves---the volume removed per acre by age for each ecoclass, prescription, and management intensity---come from the 1990 Final Environmental Impact Statement (FEIS) Appendix B \citep{usda1990umpqua} (the FORPLAN analysis volume: managed yield tables generated with the DFSIM Douglas-fir simulator, with operational-falldown adjustments, and the mountain-hemlock curve from Johnson's mountain-hemlock site equations). The economics come from the FEIS's timber-value derivations: stumpage, logging, and manufacturing costs by species (its Table B-65), the Douglas-fir price-diameter and pond-value relationships (Table B-66), and the per-ecoclass access (road) costs that make the high-elevation mountain-hemlock (CM-CE) prescriptions negatively valued. The yield-curve shapes are consistent with the LRMP's documented 95\% CMAI culmination ages by ecoclass (its Table IV-3: 85 years for CH-CW, 120 for CD-CP, 115 for CR-CF, 175 for CM-CE), and the FEIS Stage II financial analysis corroborates the thesis's anchor values (its mature-timber NPV range of \$1,800--\$7,700 per acre brackets the thesis's Table 5.4 mature-type NPVs; its managed-prescription range of \$27--\$290 per acre and the negative mountain-hemlock values match the thesis's Table 5.3 ordering and signs).

All source documents were obtained and machine-read from HathiTrust (public domain): the 1990 LRMP, the 1990 FEIS main volume and its Appendix B (the FORPLAN analysis volume), and the 1987 Draft EIS \citep{usda1987umpqua} that Daugherty used directly (the thesis's exact data vintage). The 1987 DEIS and 1990 FEIS yield tables agree, so the assembled case study uses Daugherty's exact data. The thesis's printed anchor tables (Tables 5.3 and 5.4) are reproduced as a validation check on the assembled case study. All data are typed, provenance-stamped records in the repository; the extraction of the FEIS Appendix B tables is a one-time, documented, reproducible step, and the package does not depend on the scanned source at runtime.

\hypertarget{the-sequential-replanning-simulator}{%
\subsection{The sequential-replanning simulator}\label{the-sequential-replanning-simulator}}

Dynamic inconsistency cannot be observed in a single open-loop solve; it is defined by what happens when the plan meets sequential re-optimization. Following the thesis's iterative LP simulation, we implement a sequential-replanning simulator:

\begin{enumerate}
\def\labelenumi{\arabic{enumi}.}
\tightlist
\item
  Solve the open-loop LP from the current forest state over the horizon.
\item
  Apply the current period's decision from that solution.
\item
  Advance the forest state by one period: the realized area-by-age distribution is extracted from the model and becomes the initial condition of the next subproblem.
\item
  Re-solve the open-loop LP from the realized state over the (rolling) horizon.
\item
  Repeat to the end of the simulation, recording both the plan projected at each step and the decisions actually taken.
\end{enumerate}

The inconsistency measure is the divergence between the open-loop plan's projected per-period harvest and the realized replanned trajectory---that is, the changes in decisions and in projected output levels over time between what was announced and what the re-solving planner actually does. The simulator supports a rolling fixed horizon (the default, which avoids the terminal artifacts of a shrinking horizon) and a shrinking-horizon variant; all seed-dependent runs are fixed-seeded and bit-stable.

\hypertarget{experiment-design}{%
\subsection{Experiment design}\label{experiment-design}}

The experiment grid reproduces the thesis's three experimental factors: initial forest condition (landbase), harvest policy (the even-flow tolerance), and the interest rate. The reproduced grid spans the four structured landbases of the reconstruction (landbase 1, all mature; landbases 2, 9, and 10), four discount rates (0\%, 2\%, 4\%, and 6\%), and even-flow tolerances from 1\% to 15\%. Each cell of the grid is a full sequential-replanning simulation as described above, scored for the occurrence and magnitude of dynamic inconsistency.

\hypertarget{implementation-and-reproducibility}{%
\subsection{Implementation and reproducibility}\label{implementation-and-reproducibility}}

The complete stack---typed, provenance-recorded case-study data; the \texttt{ws3} forest model; the explicit LP builder; the sequential-replanning simulator; the experiment runner; and the validation records---is openly available (see Software availability). The package is test-backed, the environment is locked, and the simulation records needed to regenerate every number reported below are produced by the repository's experiment entry points.
